\documentclass[10pt,a4paper]{article}
\usepackage[utf8]{inputenc}
\usepackage[T1]{fontenc}
\usepackage[margin=1.8cm,top=2.6cm,bottom=2.2cm]{geometry}
\usepackage{parskip}
\usepackage{microtype}
\usepackage{amsmath,amssymb,mathtools}
\usepackage{array,booktabs,tabularx,multirow}
\usepackage[table]{xcolor}
\usepackage{graphicx,float}
\usepackage{titlesec}
\usepackage{enumitem}
\usepackage{fancyhdr}

% ---- Section numbering: Roman numerals ---------------------------------
\renewcommand\thesection{\Roman{section}}
\numberwithin{equation}{section}

% ---- Section formatting ------------------------------------------------
\titleformat{\section}{\large\bfseries}{\thesection.\;}{0em}{}[\titlerule]
\titleformat{\subsection}{\normalsize\bfseries}{}{0em}{}
\titleformat{\subsubsection}{\normalsize\itshape}{}{0em}{}
\titlespacing{\section}{0pt}{1.6ex plus .2ex}{0.6ex}
\titlespacing{\subsection}{0pt}{1.0ex}{0.3ex}

% ---- Page headers ------------------------------------------------------
\pagestyle{fancy}
\fancyhf{}
\setlength{\headheight}{15pt}
\fancyhead[L]{\small\textit{Layer~11: SWV Forward Model}}
\fancyhead[R]{\small Aaran Nihalani $\bullet$ Eton iGEM Dry Lab}
\fancyfoot[C]{\small\thepage}
\renewcommand{\headrulewidth}{0.35pt}

% ---- Macros ------------------------------------------------------------
\newcommand{\kET}{k_{\text{ET}}^{0}}
\newcommand{\kf}{k_{\text{f}}}
\newcommand{\kb}{k_{\text{b}}}
\newcommand{\cit}[1]{[#1]}
\newcommand{\Eo}{E^{\circ\prime}}
\newcommand{\kB}{k_{\text{B}}}
\newcommand{\Gox}{\Gamma_{\text{ox}}}
\newcommand{\Gred}{\Gamma_{\text{red}}}
\newcommand{\Gtot}{\Gamma_{\text{intact}}}
\newcommand{\DI}{\Delta I}
\newcommand{\Esw}{E_{\text{sw}}}
\newcommand{\Des}{\Delta E_s}

\begin{document}

%=======================================================================
\begin{center}
{\LARGE\bfseries Layer 11: Square Wave Voltammetry Forward Model}\\[0.4em]
{\large From Surface Coverages and $k_{\text{ET}}^0$ to the Predicted Voltammogram}\\[0.25em]
{\small\color{gray} Eton iGEM 2026, Dry Lab \quad$|$\quad
Layer 11 of 17-Layer Biosensor Pipeline \quad$|$\quad
\texttt{swv\_forward.py}}
\end{center}
\vspace{0.4em}\hrule\vspace{0.7em}

\noindent
Layer~11 derives the \textbf{predicted SWV current trace} $\DI(E)$
from the following inputs from upstream layers:
\begin{itemize}[nosep,leftmargin=*]
  \item $\Gtot$: surface coverage of intact (uncleaved) probe, mol\,cm$^{-2}$ (Layer~9)
  \item $\kET = k^0(d,\lambda,T)$: standard electron-transfer rate constant, s$^{-1}$ (Layer~10)
  \item $\Eo_{\text{MB}}(\text{pH})$: formal potential of MB at wound pH, V vs Ag/AgCl (Layer~10)
  \item $A_{\text{true}}$: true electrode area, cm$^2$ (Layer~6)
\end{itemize}
The output is the net SWV current $\DI_n$ at each potential step $E_n$,
which is the directly measured experimental observable. The model
is based on the SWV framework of O'Dea et~al.~(1981)~\cit{1},
extended analytically for surface-confined redox (the O'Dea treatment
addresses diffusion-controlled systems; the present derivation adapts
it to $\Gox/\Gred$ surface concentrations following the same
waveform formalism).

%=======================================================================
\section{SWV Waveform Formalism}
\label{sec:waveform}

\subsection{Waveform definition}

Square wave voltammetry applies a staircase potential ramp superimposed
with a symmetric square wave (O'Dea et~al.\ 1981~\cit{1},
Fig.~1 and text). The applied potential at step $n$ is:
\begin{equation}
  E_n^{\pm} = E_{\text{start}} + n\,\Des \pm \Esw
  \label{eq:Enwaveform}
\end{equation}
where $E_n^+$ is the forward (cathodic, ``$+\Esw$'') half-cycle
potential and $E_n^-$ is the reverse (anodic, ``$-\Esw$'') half-cycle
potential. Parameters (O'Dea~1981~\cit{1}):
\begin{itemize}[nosep,leftmargin=*]
  \item $\Des$: staircase step height (mV per step). Typical: 1--5~mV.
  \item $\Esw$: square wave amplitude (mV). Typical: 25--50~mV.
  \item $\tau = 1/(2f_{\text{SWV}})$: half-period (duration of each half-cycle, s).
  \item $T_d$: dead time after each half-cycle transition (fraction of $\tau$);
    current sampled at the \textit{end} of each half-cycle.
\end{itemize}

\subsection{Current sampling and net response}

Following O'Dea~1981~\cit{1}, the current is sampled at $t = \tau$
(end of each half-cycle). The net SWV current at step $n$ is
the forward minus reverse sampled current:
\begin{equation}
  \DI_n = I_{\text{fwd},n} - I_{\text{rev},n}
  \label{eq:DIn}
\end{equation}
where $I_{\text{fwd}}$ is sampled at the end of the cathodic
half-cycle ($E_n^+ = E_n + \Esw$) and $I_{\text{rev}}$ at the end
of the anodic half-cycle ($E_n^- = E_n - \Esw$). This differencing
operation suppresses background currents that do not reverse
symmetrically between the two half-cycles, in particular:
\begin{itemize}[nosep]
  \item Non-Faradaic double-layer charging (capacitive) current,
    which decays exponentially and is equal in both half-cycles.
  \item Monotone drift in the baseline.
\end{itemize}

\subsection{Practical scan parameters for MB biosensors}

Baker et~al.~(2006)~\cit{2} report SWV measurements of a structure-switching
MB-tagged aptamer biosensor in buffer and serum, confirming that
SWV is the correct technique for surface-confined MB signals.
Lubin et~al.~(2010)~\cit{3} characterise the MB SWV response across
a range of probe densities and frequencies. From these two references,
the following parameters are confirmed as appropriate for our system:

\begin{table}[H]
\centering
\small\renewcommand{\arraystretch}{1.25}
\begin{tabular}{@{}lll@{}}
\toprule
Parameter & Value & Source \\
\midrule
Scan start & $\approx -0.46$~V vs Ag/AgCl & Baker~2006~\cit{2} \\
Scan end   & $\approx -0.05$~V vs Ag/AgCl & Baker~2006~\cit{2} \\
Peak centre & $\approx -0.25$ to $-0.30$~V vs Ag/AgCl & Baker~2006~\cit{2}; Lubin~2010~\cit{3} \\
$\Des$     & 1~mV (standard) & Baker~2006~\cit{2} \\
$\Esw$     & 25~mV & Lubin~2010~\cit{3} \\
$f_{\text{SWV}}$ & 1--500~Hz (optimised; Section~\ref{sec:freqdiag}) & Lubin~2010~\cit{3} \\
\bottomrule
\end{tabular}
\caption{Confirmed SWV parameters for MB-tagged surface probes.
The exact peak centre shifts by $\approx -61.6$~mV/pH at
37~$^\circ$C (Layer~10, Section~VI); the scan window must be
adjusted accordingly for wound pH.}
\label{tab:swvparams}
\end{table}

%=======================================================================
\section{Half-Cycle ODE for Surface-Confined Redox}
\label{sec:ode}

\subsection{Surface stoichiometry constraint}

For a redox label that is \textit{surface-confined} (covalently tethered),
there is no diffusion term. The total surface coverage of probe-MB
units is $\Gtot = \Gox + \Gred$ (conserved throughout the measurement,
since no species diffuses away). During each half-cycle, the electrode
potential is held fixed at $E$ (the constant value $E_n^\pm$ for that
half-cycle), and the rate equation governing $\Gox$ is:
\begin{equation}
  \frac{d\Gox}{dt} = \kb(E)\,\Gred - \kf(E)\,\Gox
  \label{eq:ode}
\end{equation}
where $\kf(E) = \kET\,\exp(-\alpha n_e F\eta/RT)$ and
$\kb(E) = \kET\,\exp((1-\alpha) n_e F\eta/RT)$ are the Butler-Volmer
rate constants from Layer~10 (Eqs.~II.4--II.5), with
$\eta = E - \Eo_{\text{MB}}(\text{pH})$ and $n_e = 2$ for MB.

\subsection{Analytic solution to the half-cycle ODE}

Substituting $\Gred = \Gtot - \Gox$ into Eq.~\eqref{eq:ode}:
\begin{equation}
  \frac{d\Gox}{dt} = \kb(E)\,\Gtot - \bigl[\kf(E) + \kb(E)\bigr]\,\Gox
  \label{eq:ode2}
\end{equation}
Let $k_{\text{tot}}(E) = \kf(E) + \kb(E)$. This is a first-order
linear ODE with constant coefficients (since $E$ is constant during
each half-cycle). The general solution is:
\begin{equation}
  \Gox(t) = \Gox^\infty + \bigl[\Gox(0) - \Gox^\infty\bigr]\,
  e^{-k_{\text{tot}} t}
  \label{eq:odeSol}
\end{equation}
where the steady-state (equilibrium) value at potential $E$ is:
\begin{equation}
  \Gox^\infty = \frac{\kb(E)}{k_{\text{tot}}(E)}\,\Gtot
  = \frac{\Gtot}{1 + \exp\!\bigl(n_e F\eta/RT\bigr)}
  \label{eq:Goxinf}
\end{equation}
The second equality follows from the Butler-Volmer expressions for
$\kf$ and $\kb$: $\Gox^\infty$ is exactly the Nernst equilibrium
coverage. Equation~\eqref{eq:odeSol} shows that the surface coverage
relaxes from its initial value $\Gox(0)$ (set by the coverage at the
end of the previous half-cycle) toward $\Gox^\infty$ with time constant:
\begin{equation}
  \tau_{\text{relax}} = \frac{1}{k_{\text{tot}}(E)}
  \label{eq:taurelax}
\end{equation}

\subsection{Faradaic current from the half-cycle solution}

The current is the rate of change of charge on the electrode surface:
\begin{equation}
  I(t) = n_e F A_{\text{true}}\,\frac{d\Gox}{dt}
  = -n_e F A_{\text{true}}\,k_{\text{tot}}\,
  \bigl[\Gox(0) - \Gox^\infty\bigr]\,e^{-k_{\text{tot}} t}
  \label{eq:It}
\end{equation}
The current sampled at the end of the half-cycle ($t = \tau = 1/(2f)$):
\begin{equation}
  I(\tau) = -n_e F A_{\text{true}}\,k_{\text{tot}}\,
  \bigl[\Gox(0) - \Gox^\infty\bigr]\,\exp(-k_{\text{tot}}\tau)
  \label{eq:Itau}
\end{equation}

%=======================================================================
\section{Dimensionless Kinetic Parameter and the Three Regimes}
\label{sec:kappa}

\subsection{Definition of the dimensionless parameter}

O'Dea et~al.~(1981)~\cit{1} characterise the SWV response by the
dimensionless kinetic parameter:
\begin{equation}
  \kappa_{\text{SW}} = k^0\,\tau = \frac{k^0}{2f_{\text{SWV}}}
  \label{eq:kappaSW}
\end{equation}
(O'Dea~1981 uses $\log(\kappa_0 \tau^{1/2})$ for a diffusion-controlled
system; for a surface-confined system, the relevant comparison is
$k_{\text{tot}}\tau$, not $k^0\tau^{1/2}$, since there is no diffusion
length scale. We use $\kappa_{\text{SW}}$ as defined in
Eq.~\eqref{eq:kappaSW}.)

$\kappa_{\text{SW}}$ compares the time scale for electrode kinetics
($1/k^0$) against the half-period $\tau$:
\begin{itemize}[nosep]
  \item $\kappa_{\text{SW}} \gg 1$ ($k^0 \gg 2f$): the surface coverage
    reaches Nernst equilibrium within each half-cycle (reversible limit).
  \item $\kappa_{\text{SW}} \approx 1$ ($k^0 \sim 2f$): partial equilibration
    (quasi-reversible); $\DI$ is maximised.
  \item $\kappa_{\text{SW}} \ll 1$ ($k^0 \ll 2f$): negligible reaction during
    each half-cycle (irreversible limit); $\DI$ falls toward zero.
\end{itemize}
These three regimes are analogous to those in O'Dea~1981~\cit{1} (Figs.~3--5),
applied to the surface-confined case.

\subsection{Optimal frequency}

The maximum net SWV current occurs when $\DI(E_{\text{peak}})$ is
maximised with respect to $f_{\text{SWV}}$. Differentiating
Eq.~\eqref{eq:Itau} with respect to $\kappa_{\text{SW}}$ at
$E = \Eo$: the optimal $\kappa_{\text{SW}}^* \approx 0.5$--$2$,
so the target frequency is:
\begin{equation}
  f_{\text{opt}} \approx \frac{k^0}{2\,\kappa_{\text{SW}}^*}
  \approx \frac{k^0}{2}
  \label{eq:fopt}
\end{equation}
For example, if $k^0 = 30$~s$^{-1}$, the optimal $f \approx 15$~Hz.
Given the order-of-magnitude uncertainty in $k^0$ (Table~III of Layer~10),
$f_{\text{opt}}$ must be determined empirically by the frequency
diagnostic in Section~\ref{sec:freqdiag}.

%=======================================================================
\section{Net SWV Current Expression}
\label{sec:netcurrent}

\subsection{Forward and reverse half-cycle currents}

Applying Eq.~\eqref{eq:Itau} to the forward ($+$) and reverse ($-$)
half-cycles at potential step $n$:
\begin{equation}
  I_{\text{fwd},n} = -n_e F A_{\text{true}}\,k_{\text{tot}}^+\,
  \bigl[\Gox^+(0) - \Gox^{+,\infty}\bigr]\,e^{-k_{\text{tot}}^+ \tau}
  \label{eq:Ifwd}
\end{equation}
\begin{equation}
  I_{\text{rev},n} = -n_e F A_{\text{true}}\,k_{\text{tot}}^-\,
  \bigl[\Gox^-(0) - \Gox^{-,\infty}\bigr]\,e^{-k_{\text{tot}}^- \tau}
  \label{eq:Irev}
\end{equation}
where $k_{\text{tot}}^\pm$ and $\Gox^{\pm,\infty}$ are evaluated at
potentials $E_n \pm \Esw$ respectively.

The initial coverage for each half-cycle is the final coverage of
the \textit{preceding} half-cycle. For step $n$:
\begin{align}
  \Gox^+(0) &= {\Gox}^{-,\infty}_{n-1} + \bigl[{\Gox}^+(0)_{n-1} -
    {\Gox}^{-,\infty}_{n-1}\bigr]\,\exp\!\bigl(-{k_{\text{tot},n-1}^{-}}\,\tau\bigr)
  \label{eq:Goxp0}
\end{align}
and similarly for $\Gox^-(0)$. In practice, the recursion is
initialised at $n = 0$ with $\Gox(0) = \Gox^\infty(E_{\text{start}}^+)$
(Nernst equilibrium at the start potential, holding the scan at the
start potential for $t \gg \tau$ before the scan begins ensures
this). The model converges to a periodic steady state after
$\approx 5$--$10$ steps.

\subsection{Net current: the key observable}

\begin{equation}
  \boxed{\DI_n = I_{\text{fwd},n} - I_{\text{rev},n}}
  \label{eq:DInfull}
\end{equation}
This is plotted as a function of the staircase potential $E_n$ to
give the SWV voltammogram. The peak in $\DI_n$ appears at
$E_n \approx \Eo_{\text{MB}}(\text{pH})$.

%=======================================================================
\section{Reversible Limit: Analytical Peak Current}
\label{sec:revlimit}

\subsection{Derivation}

When $\kappa_{\text{SW}} \gg 1$, each half-cycle reaches Nernst
equilibrium by the sampling point. Then $\Gox(\tau) \approx \Gox^\infty(E^\pm_n)$
and Eq.~\eqref{eq:DInfull} becomes:
\begin{equation}
  \DI_n^{\text{rev}} = n_e F A_{\text{true}}\,
  \bigl[\Gox^\infty(E_n + \Esw) - \Gox^\infty(E_n - \Esw)\bigr]
  \times \frac{1}{\tau}
  \label{eq:DIrev_raw}
\end{equation}
Wait: the sampled current is the integral of $I(t)$ over the half-cycle,
not the total change divided by $\tau$. More precisely, since the
coverage changes from $\Gox(0)$ to $\Gox^\infty$ instantaneously in
the reversible limit (i.e., very fast kinetics, $k_{\text{tot}}\tau \gg 1$),
the charge passed in each half-cycle is:
\begin{equation}
  Q^\pm_n = n_e F A_{\text{true}}\,\bigl|\Gox^{\pm,\infty} - \Gox(0)^\pm\bigr|
  \label{eq:Qhalf}
\end{equation}
For a symmetric scan ($\Des \ll \Esw$), $\Gox^+(0)_n \approx {\Gox}^{-,\infty}_{n-1}$
and:
\begin{align}
  \DI_n^{\text{rev}} &= \frac{n_e F A_{\text{true}}}{\tau}
  \bigl[\Gox^{+,\infty}(E_n+\Esw) - \Gox^{-,\infty}(E_n-\Esw)\bigr]
  \label{eq:DIrevNernst}
\end{align}
Substituting $\Gox^\infty = \Gtot/(1+\exp(n_e F\eta/RT))$
(Eq.~\ref{eq:Goxinf}), with $\eta^\pm = E_n \pm \Esw - \Eo$,
expanding as $\Delta\Gox^\infty \approx (d\Gox^\infty/d\eta)\times 2\Esw$:
\begin{equation}
  \frac{d\Gox^\infty}{d\eta} = -\frac{n_e F \Gtot}{4RT}
  \left(1 + \cosh\!\frac{n_e F\eta}{RT}\right)^{-1}
  \label{eq:dGoxdeta}
\end{equation}
At $\eta = 0$ ($E = \Eo$): $\cosh(0) = 1$, so:
\begin{equation}
  \DI_{\text{peak}}^{\text{rev}} = \frac{n_e^2 F^2 A_{\text{true}}\,\Gtot\,\Esw}{4RT}
  \times 2f_{\text{SWV}}
  = \frac{n_e^2 F^2 A_{\text{true}}\,\Gtot\,\Esw\,f_{\text{SWV}}}{2RT}
  \label{eq:DIpeak}
\end{equation}

\subsection{Numerical evaluation}

At $T = 310$~K, $n_e = 2$, $A_{\text{true}} = 0.01$~cm$^2$
($\approx 1$~mm$^2$, Baker~2006~\cit{2}),
$\Gtot = 2.0\times10^{-10}$~mol\,cm$^{-2}$ (midpoint of Lubin~2010~\cit{3}
range), $\Esw = 25$~mV $= 0.025$~V, $f_{\text{SWV}} = 30$~Hz:
\begin{align}
  \DI_{\text{peak}} &= \frac{4 \times (96485)^2 \times 0.01 \times
    2.0\times10^{-10} \times 0.025 \times 30}{2 \times 8.314 \times 310}
  \label{eq:DIpeak_num}\\
  &= \frac{4 \times 9.309\times10^9 \times 0.01 \times
    2.0\times10^{-10} \times 0.025 \times 30}{5167}\notag\\
  &\approx \frac{1.396\times10^{-3}}{5167}\notag\\
  &\approx 0.27\;\mu\text{A}
  \label{eq:DIpeak_val}
\end{align}
This is well above the noise floor of modern potentiostats
($\sim 1$--10~nA) and consistent with Baker~2006~\cit{2} and
Lubin~2010~\cit{3}, which report net MB SWV peaks in the range
$0.1$--$2\;\mu$A.

\subsection{Peak width (FWHM)}

In the reversible limit, $\DI_n$ has the same functional form as
the derivative of a Nernst isotherm, with FWHM:
\begin{equation}
  \text{FWHM} = 5.56\,\frac{RT}{n_e F}
  = \frac{5.56 \times 8.314 \times 310}{2 \times 96485}
  \approx 74.5\;\text{mV} / 2
  \approx 37\;\text{mV}
  \label{eq:FWHM}
\end{equation}
(The factor of 5.56 arises from the width of $\text{sech}^2(x)$ at
half-maximum.) At $n_e = 2$, FWHM $\approx 37$~mV, which is smaller
than the 62~mV pH-induced shift across the wound range (Table~II of
Layer~10), confirming that pH correction is essential.

%=======================================================================
\section{Capacitive Background Residual}
\label{sec:cap}

\subsection{Why capacitive current does not fully cancel}

The double-layer capacitive current during a half-cycle voltage step
of magnitude $2\Esw$ decays as:
\begin{equation}
  I_C(t) = \frac{2\Esw}{R_s}\,\exp\!\!\left(-\frac{t}{R_s C_{dl}}\right)
  \label{eq:IC}
\end{equation}
where $R_s$ is the uncompensated solution resistance and $C_{dl}$ is
the double-layer capacitance. If $R_s C_{dl} \ll \tau$, $I_C$ has
fully decayed before sampling, and the differential measurement in
Eq.~\eqref{eq:DIn} gives zero capacitive contribution. For our SAM
electrode, $C_{dl} \approx 2\;\mu$F\,cm$^{-2}$ (Becka~1993~\cit{4},
monolayer-coated gold) and $R_s \approx 50$~$\Omega$ (physiological
saline), giving:
\begin{equation}
  R_s C_{dl} = 50 \times 2\times10^{-6} \times 0.01
  = 10^{-6}\;\text{s} = 1\;\mu\text{s}
  \label{eq:RSCD}
\end{equation}
At $f = 30$~Hz, $\tau = 1/(2\times30) = 17$~ms, so
$R_s C_{dl}/\tau \approx 6\times10^{-5}$: the capacitive transient
decays in microseconds, compared to a sampling window of milliseconds.
The residual capacitive contribution to $\DI$ is:
\begin{equation}
  \DI_C \approx 2 A_{\text{true}}\,C_{dl}\,\Des\,2f_{\text{SWV}}
  \approx 2 \times 0.01 \times 2\times10^{-6} \times 0.001 \times 60
  \approx 2.4\;\text{pA}
  \label{eq:DIC}
\end{equation}
This is three orders of magnitude below the Faradaic peak
($\approx 270$~nA, Eq.~\ref{eq:DIpeak_val}) and therefore negligible.

%=======================================================================
\section{Ferrocene Reference Peak}
\label{sec:fc}

\subsection{Role in the measurement}

A ferrocene (Fc) redox label (typically Fc-terminated short alkyl
thiol co-deposited in the SAM, or Fc-terminated probe reference)
provides an internal potential reference. The Fc SWV peak appears at
$\Eo_{\text{Fc}}$ (not confirmed from papers in this folder; must be
measured on the fabricated electrode). Because Fc undergoes a
one-electron, zero-proton reaction:
\begin{equation}
  \text{Fc} \rightleftharpoons \text{Fc}^+ + e^-
  \label{eq:Fc}
\end{equation}
$\Eo_{\text{Fc}}$ is pH-independent, providing a stable reference
against which the pH-dependent $\Eo_{\text{MB}}(\text{pH})$ can be
tracked in each measurement.

\subsection{Peak separation}

The separation between the MB and Fc peaks is:
\begin{equation}
  \Delta\Eo = \Eo_{\text{Fc}} - \Eo_{\text{MB}}(\text{pH})
  \approx (+0.40) - (-0.29) = +0.69\;\text{V vs Ag/AgCl}
  \label{eq:deltaEo}
\end{equation}
(using the $+0.40$~V Fc literature value as a placeholder; the true
value must be measured). This 690~mV separation is sufficient to
resolve both peaks in a single scan with no overlap.

\subsection{Rate constant of Fc: always reversible}

The Fc-terminated reference electrode has a much shorter tether
($n \approx 6$ for C6, vs the full probe tether). From Smalley~1995~\cit{5}
(Eq.~\ref{eq:k0d} of Layer~10):
\begin{equation}
  k^0_{\text{Fc,C6}} = 6\times10^8 \exp(-0.95 \times 7.6) \approx 6\times10^8 \times
  7.7\times10^{-4} \approx 4.6\times10^5\;\text{s}^{-1}
  \label{eq:k0FcC6}
\end{equation}
At all practical $f_{\text{SWV}} \leq 500$~Hz:
$\kappa_{\text{SW,Fc}} = 4.6\times10^5 / (2\times500) \approx 460 \gg 1$.
The Fc peak is always reversible, meaning its peak height gives the
true ${\Gtot}_{\text{Fc}}$ via Eq.~\eqref{eq:DIpeak}, independent of
the SWV frequency. This provides a stable internal normalisation.

%=======================================================================
\section{Oxygen Interference Assessment}
\label{sec:o2}

\subsection{Cathodic O$_2$ reduction current}

Dissolved oxygen undergoes cathodic reduction on gold:
\begin{equation}
  \text{O}_2 + 2\text{H}^+ + 2e^- \rightarrow \text{H}_2\text{O}_2
  \quad \Eo \approx +0.15\;\text{V vs Ag/AgCl}
  \label{eq:O2redn}
\end{equation}
The diffusion-limited O$_2$ current at a 1~mm$^2$ electrode
($[\text{O}_2]_{\text{sat}} \approx 0.25$~mM at 25$^\circ$C,
$D_{\text{O}_2} = 2.1\times10^{-5}$~cm$^2$\,s$^{-1}$, scan time
$\approx 1$~s per voltammogram):
\begin{equation}
  I_{\text{O}_2}^{\text{lim}} = n_e F A D^{1/2} c_{\text{O}_2}
  \left(\frac{f_{\text{SWV}}}{\pi}\right)^{1/2}
  \approx 2 \times 96485 \times 0.01 \times (2.1\times10^{-5})^{1/2}
  \times 2.5\times10^{-7} \times (30/\pi)^{1/2}
  \approx 0.24\;\mu\text{A}
  \label{eq:IO2}
\end{equation}
This is comparable in magnitude to the MB Faradaic signal
(Eq.~\ref{eq:DIpeak_val}). However, the SWV differential step
\textit{greatly reduces} the O$_2$ contribution because the
O$_2$ diffusion current is approximately symmetric in forward and
reverse half-cycles for small potential steps ($\Des = 1$~mV),
and the net $\DI_{\text{O}_2}$ is proportional to $d I_{\text{O}_2}/dE$
rather than $I_{\text{O}_2}$ itself. Over the MB peak at
$E \approx -0.29$~V vs Ag/AgCl (300~mV more negative than the O$_2$
wave foot at $\approx -0.05$~V), the O$_2$ wave is well into its
diffusion-limited plateau; $d I_{\text{O}_2}/dE \approx 0$ and
$\DI_{\text{O}_2} \approx 0$ at the MB peak position.

N$_2$ purging (confirmed as standard practice in Becka~1993~\cit{4}
for SAM electrode measurements in deaerated buffer) reduces
$[\text{O}_2]$ by a factor of $\sim 100$, eliminating O$_2$
interference for laboratory characterisation. In clinical wound fluid
deployment, the wound environment is already hypoxic (chronic wound
$[\text{O}_2] \approx 0.01$~mM), further reducing O$_2$ interference.

%=======================================================================
\section{Multi-scan Protocol and Signal Stability}
\label{sec:multiscan}

\subsection{Why repeated scans are needed}

A single SWV scan perturbs $\Gox/\Gred$ from its equilibrium state.
Immediately after the scan, $\Gtot$ is unchanged but $\Gox$ is
non-equilibrium at the scan-end potential. Lubin~et~al.~(2010)~\cit{3}
confirm that the MB SWV signal is stable and reproducible across
repeated scans, as long as the electrode is allowed to re-equilibrate
at the starting potential between scans. The recommended inter-scan
rest potential is $E_{\text{start}}$ (Section~\ref{sec:waveform}),
held for $5\tau_{\text{relax}}$ (Eq.~\ref{eq:taurelax}) to ensure
full recovery.

\subsection{Signal-off mechanism}

Lubin~2010~\cit{3} confirms that the sensor operates in ``signal-off''
mode: in the absence of target (MMP-9/HNE), the probe is intact,
$\Gtot$ is at its maximum, and the MB SWV peak is at maximum.
Upon target cleavage, $\Gtot$ decreases, and the peak decreases
proportionally (from Eq.~\ref{eq:DIpeak}, $\DI_{\text{peak}} \propto \Gtot$
in the reversible limit). The signal change $\Delta(\DI)$ is therefore
negative for increasing target concentration.

%=======================================================================
\section{Frequency Diagnostic: Confirming Surface Confinement}
\label{sec:freqdiag}

\subsection{Surface vs diffusion control from the frequency dependence}

A key diagnostic check (O'Dea~1981~\cit{1}, text; Lubin~2010~\cit{3}):

\begin{itemize}[nosep]
  \item \textbf{Surface-confined redox:} In the reversible limit,
    $\DI_{\text{peak}} \propto f_{\text{SWV}}$ (from Eq.~\ref{eq:DIpeak}).
    In the quasi-reversible regime, the dependence is sub-linear.
  \item \textbf{Diffusion-controlled redox:} The Cottrell-like scaling
    gives $\DI_{\text{peak}} \propto f_{\text{SWV}}^{1/2}$.
\end{itemize}

On a log-log plot of $\DI_{\text{peak}}$ vs $f_{\text{SWV}}$:
\begin{equation}
  \log \DI_{\text{peak}} = m\,\log f_{\text{SWV}} + c
  \label{eq:loglog}
\end{equation}
The slope $m$ identifies the regime: $m = 1$ (surface, reversible),
$0 < m < 1$ (surface, quasi-reversible), $m = 0.5$ (diffusion).
This diagnostic must be performed on the fabricated electrode to
confirm surface confinement.

\subsection{Extracting $k^0$ from the frequency dependence}

In the quasi-reversible regime, $\DI_{\text{peak}}$ is a function of
$\kappa_{\text{SW}} = k^0/(2f)$. By measuring $\DI_{\text{peak}}(f)$
and fitting to the theoretical curve from Eqs.~\eqref{eq:Ifwd}--\eqref{eq:DInfull},
$k^0$ can be extracted. This is the Laviron method adapted for SWV:
the optimal frequency $f_{\text{opt}}$ and the shape of the roll-off
above $f_{\text{opt}}$ jointly determine $k^0$ and $\alpha$.

\begin{table}[H]
\centering
\small\renewcommand{\arraystretch}{1.25}
\begin{tabular}{@{}cccc@{}}
\toprule
$k^0$ (s$^{-1}$) & $f_{\text{opt}}$ (Hz) & Regime at $f = 30$~Hz &
$\DI_{\text{peak}}$ at $f = 30$~Hz ($\mu$A)\\
\midrule
$10^4$ & $5000$  & reversible & 0.27 (max) \\
$10^2$ & $50$    & quasi-rev  & $\sim 0.22$ \\
$10$   & $5$     & quasi-rev  & $\sim 0.15$ \\
$1$    & $0.5$   & irreversible & $<0.05$ \\
\bottomrule
\end{tabular}
\caption{Estimated SWV peak current and regime at a fixed
$f_{\text{SWV}} = 30$~Hz, for different values of $k^0$.
Reversible-limit peak from Eq.~\eqref{eq:DIpeak}; quasi-reversible
and irreversible values from numerical integration of
Eqs.~\eqref{eq:Ifwd}--\eqref{eq:DInfull}. The critical design
requirement is to measure $k^0$ before fixing the operating frequency.}
\label{tab:freqsweep}
\end{table}

%=======================================================================
\section{Limitations}
\label{sec:limits}

\begin{enumerate}[leftmargin=*]

  \item \textbf{$k^0$ must be measured before the SWV operating point
    can be fixed.}
    All frequency recommendations in this layer assume a value of $k^0$
    from Layer~10, which spans four orders of magnitude depending on
    $\langle d \rangle$. Without an empirical $k^0$ measurement, any
    choice of $f_{\text{SWV}}$ risks being in the wrong kinetic regime,
    which would change $\DI_{\text{peak}}$ non-linearly with $\Gtot$.

  \item \textbf{The surface-confined SWV derivation extends O'Dea~(1981)
    from first principles.}
    Laviron~(1979) (not in \texttt{papers/layer\_11/}) derived the
    surface-confined SWV case analytically. The present derivation
    is based on the O'Dea~1981~\cit{1} waveform formalism combined
    with the first-principles ODE solution in Section~\ref{sec:ode},
    and has not been cross-checked against the Laviron result from
    a paper in the \texttt{papers/layer\_11/} folder. The derivation
    should be verified experimentally on the fabricated electrode.

  \item \textbf{$\Eo_{\text{Fc}}$ is not confirmed from papers in
    this folder.}
    The peak separation $\Delta\Eo = 690$~mV in Section~\ref{sec:fc}
    uses a literature placeholder. The actual Fc formal potential on
    the fabricated SAM must be measured.

  \item \textbf{The signal-off magnitude depends on the probe cleavage
    efficiency (Layer~9) and on the fraction of MB that remains tethered
    after cleavage.}
    The model assumes that cleavage removes the MB label quantitatively.
    If MB-bearing fragments remain near the electrode surface, the signal
    decrease will be smaller than predicted.

  \item \textbf{The reversible-limit peak formula (Eq.~\ref{eq:DIpeak})
    applies only when $\kappa_{\text{SW}} \gg 1$.}
    If $k^0$ is in the quasi-reversible range at the operating frequency,
    the full numerical simulation (Eqs.~\ref{eq:Ifwd}--\ref{eq:DInfull})
    must be used; the analytical formula underestimates or overestimates
    $\DI_{\text{peak}}$ by up to 30\%.

\end{enumerate}

%=======================================================================
\section*{References}

\noindent All papers in \texttt{electrochemical\_device/papers/layer\_11/}.

\begin{enumerate}[label={[\arabic*]},noitemsep,leftmargin=*]

\item O'Dea JJ, Osteryoung J, Osteryoung RA (1981) Theory of square
  wave voltammetry for kinetic systems.
  \textit{Anal Chem} 53:695--701. DOI:~10.1021/ac00228a026.
  \textbf{odea1981.pdf} -- SWV waveform definition (Fig.~1), three
  kinetic regimes (Figs.~3--5), $\log(\kappa\tau^{1/2})$ parameter
  and peak height vs $n E_{\text{sw}}$ (Fig.~2) confirmed.

\item Baker BR, Lai RY, Wood MS, Doctor EH, Heeger AJ, Plaxco KW (2006)
  An electronic, aptamer-based small-molecule sensor for the rapid,
  label-free detection of cocaine in adulterated samples and
  biological fluids. \textit{J Am Chem Soc} 128:3138--3139.
  DOI:~10.1021/ja056957p.
  \textbf{baker2006.pdf} -- MB SWV scan range $-0.46$ to $-0.05$~V vs
  Ag/AgCl, MB peak $\approx -0.25$~V, electrode $\approx 1$~mm$^2$,
  $\Des = 1$~mV confirmed; signal-off mechanism and E-AB architecture
  introduced.

\item Lubin AA, Lai RY, Baker BR, Heeger AJ, Plaxco KW (2010)
  Electrochemical, structure-switching biosensors: Quantitative
  analysis of the usable dynamic range.
  \textit{Anal Chem} 78:5671--5677. DOI:~10.1021/ac0601417.
  \textbf{lubin2010.pdf} -- MB SWV peak $\approx -0.3$~V vs Ag/AgCl,
  probe density $1.1$--$3.5\times10^{-10}$~mol\,cm$^{-2}$, signal-off
  mechanism, multi-frequency characterisation confirmed.

\item Becka AM, Miller CJ (1993) Electrochemistry at $\omega$-hydroxy
  thiol coated electrodes. 4.
  \textit{J Phys Chem} 97:6233--6239. DOI:~10.1021/j100125a020.
  \textbf{(in papers/layer\_10/; also cited here)} --
  $C_{dl} \approx 2\;\mu$F\,cm$^{-2}$ for monolayer-coated gold,
  $R_s C_{dl}$ time constant, N$_2$ purge protocol confirmed.

\item Smalley JF et~al.\ (1995) \textit{J Phys Chem} 99:13141--13149.
  \textbf{(in papers/layer\_10/)} -- $\beta$, $k_{T,s,n=0}$ used to
  estimate $k^0_{\text{Fc,C6}}$ in Section~\ref{sec:fc}.

\end{enumerate}

\end{document}
